% Source: https://artofproblemsolving.com/community/c3h2780957_help_on_number_theory_question_mathcounts
% Stable identity: aops-topic-2780957-q01
Group the powers of $2$ by their number of digits. Doubling a number increases its digit count exactly when the number starts with $5,6,7,8$ or $9$, and then the result starts with $1$. So each group starts with a power whose leading digit is $1$, and within a group the leading digits follow one of the patterns
\[1,2,4,8\qquad1,2,4,9\qquad1,2,5\qquad1,3,6\qquad1,3,7.\]
(From a leading $1$ the next leading digit is $2$ or $3$; from $2$ it is $4$ or $5$; from $3$ it is $6$ or $7$; from $4$ it is $8$ or $9$; and $5,\dots,9$ end the group.)

The powers $2^0,\dots,2^{999}$ form the complete groups of $1,2,\dots,301$ digits, and $2^{1000}$ alone starts the $302$-digit group (its leading digit is $1$). If $x$ of the $301$ complete groups have four members, then $4x+3(301-x)=1000$, so $x=97$. Each four-member group contains exactly one $8$ or one $9$, and $8$ appears $52$ times, so $9$ appears $97-52=\boxed{45}$ times.
